\section{Обзор литературы}\label{sec:review}

\subsection{Результаты предыдущей НИР}

В предыдущей работе была реализована процедура сканирования двумерного
пространства $(\eta_0,\eta_1)$ для однослойной сети, введён бинарный критерий
сходимости (среднее нормализованного loss по последним 20 итерациям меньше
единицы) и построены цветовые карты обучаемости для активаций tanh, ReLU и
линейной сети. Измеренные box-counting размерности составили $1.547$ для tanh,
$1.612$ для ReLU и $1.183$ для линейной сети; последовательность из семи уровней
зума дала медианную размерность $1.548$ при стандартном отклонении $0.009$.
На игрушечной задаче квадратичной регрессии были воспроизведены пять фаз
динамики: монотонная сходимость, catapult, периодические орбиты, хаос и
расходимость. Настоящая работа наследует программную архитектуру и метод
измерения размерности, но расширяет пространство гиперпараметров и уточняет
классификацию режимов.

\subsection{Фрактальная граница обучаемости}

Явление было обнаружено и описано Sohl-Dickstein~\cite{sohl2024}. Основное
наблюдение состоит в том, что множество гиперпараметров, при которых обучение
сходится, имеет границу с самоподобной структурой, наблюдаемой на многих
порядках масштаба, а измеренные размерности лежат в диапазоне примерно от
$1.1$ до $2.0$ в зависимости от активации и выбранной пары гиперпараметров.
Ключевой аргумент состоит в аналогии: обучение --- это итерация отображения
в пространстве параметров, а классические фракталы (множество Мандельброта,
множества Жюлиа) возникают именно как границы областей ограниченности при
итерации простых отображений.

Torkamandi и соавторы~\cite{torkamandi2025} перенесли эксперимент на
трансформерные архитектуры типа decoder-only, обучаемые оптимизатором Adam, и
получили фрактальные границы с размерностями порядка $1.5$--$2.0$, что
указывает на универсальность явления и на то, что адаптивные методы сами по
себе фрактальность не устраняют. Liu, Arnould и Sohl-Dickstein~\cite{liu2024}
показали, что для возникновения фрактальной границы достаточно тривиальной
невыпуклости: добавление косинусоидального возмущения к квадратичной функции
потерь уже порождает фрактальную границу, причём размерность растёт с
«шероховатостью» возмущения. Это важно методологически: наблюдаемая
фрактальность не требует ни глубины сети, ни большого объёма данных.

\subsection{Грань устойчивости и крупные шаги обучения}

Cohen и соавторы~\cite{cohen2021} обнаружили, что при полнобатчевом обучении
максимальное собственное значение гессиана растёт до величины $2/\eta$ и далее
стабилизируется около этого порога, а loss становится немонотонным на коротких
масштабах, продолжая убывать на длинных. Этот режим получил название грани
устойчивости. Он объясняет, почему наилучшие по скорости конфигурации
располагаются вблизи границы обучаемости, а не в глубине области сходимости.
В последующей работе~\cite{cohen2022} тот же эффект был описан и для адаптивных
методов, где роль порога играет предобусловленная резкость; это делает
нетривиальным вопрос о том, сглаживает ли Adam границу.

Lewkowycz и соавторы~\cite{lewkowycz2020} описали catapult-механизм: при
достаточно крупном шаге loss сначала растёт, а затем система переходит в
область меньшей кривизны и сходимость восстанавливается. Kong и
Tao~\cite{kong2020} показали, что детерминированный градиентный спуск с крупным
шагом на многомасштабной целевой функции ведёт себя стохастически: если шаг
разрешает макромасштаб, но не микромасштаб, динамика становится хаотической и
возникают эффекты, аналогичные шуму SGD. Chen и соавторы~\cite{chen2023}
провели детальный анализ квадратичной регрессии, где увеличение шага
последовательно приводит систему через монотонную сходимость, catapult,
периодические орбиты, хаос по Ли--Йорке~\cite{liyorke1975} и расходимость.
Универсальный сценарий перехода к хаосу через каскад удвоения периода описан
Feigenbaum~\cite{feigenbaum1978}.

\subsection{Методы оптимизации}

Метод тяжёлого шарика предложен Поляком~\cite{polyak1964}; ускоренный метод
с иной схемой экстраполяции --- Нестеровым~\cite{nesterov1983}. Adam введён
Kingma и Ba~\cite{kingma2015}, RMSprop --- Tieleman и
Hinton~\cite{tieleman2012}. Loshchilov и Hutter~\cite{loshchilov2019} показали,
что для адаптивных методов существенно, добавляется ли затухание весов к
градиенту или применяется развязанно; в настоящей работе используется первый,
классический вариант, что необходимо учитывать при сопоставлении результатов.
Обзор методов оптимизации для крупномасштабного машинного обучения дан в
работе Bottou, Curtis и Nocedal~\cite{bottou2018}; связь размера батча и шага
обучения исследована Smith и соавторами~\cite{smith2018}. Систематическое
изложение архитектур, активаций и алгоритмов обучения содержится в
монографии Goodfellow, Bengio и Courville~\cite{goodfellow2016}.

\subsection{Фрактальная геометрия и динамические системы}

Математические основы фрактальной геометрии и определения размерности, включая
размерность Минковского--Булиганда, изложены у Falconer~\cite{falconer2014};
классическое изложение с акцентом на природные примеры дано
Mandelbrot~\cite{mandelbrot1982}. Теория нелинейных динамических систем,
бифуркации и переход к хаосу описаны у Strogatz~\cite{strogatz2015}.
Практические аспекты подсчёта box-counting размерности по бинарным изображениям
и типичные источники систематической ошибки рассмотрены в документации
пакета PoreSpy~\cite{porespy2019}.

\subsection{Выводы по обзору}

Литература устанавливает, что фрактальность границы обучаемости --- устойчивое
явление, наблюдаемое от простейших невыпуклых функций до трансформеров, и что
оно тесно связано с динамикой вблизи грани устойчивости. Однако
систематического сравнения сложности границы между оптимизаторами в
сопоставимых сечениях пространства гиперпараметров в доступной литературе не
обнаружено; в частности, вопрос о том, сглаживает ли покоординатная
нормализация градиентов границу, остаётся открытым. Настоящая работа
предоставляет количественные данные по этому вопросу.
